\documentclass[a4paper, 10pt, ]{article}

\input{../../KUT000/preamble_v1.0.tex}

% -----------------------------------------------------------------------------

\def\oznacenieCelku{Kolekcia učebných textov}

% -----------------------------------------------------------------------------


\def\KUTporadoveCislo{021}

\def\oznacenieVerzie{v0.97}
% \def\oznacenieVerzie{\phantom{v1.0}}

\def\mesiacRok{marec 2026}

\def\authorslabel{AS, MT}






% -----------------------------------------------------------------------------

\begin{document}

% -----------------------------------------------------------------------------
% Uvodny nadpis

\noindent
\parbox[t][18mm][c]{0.3\textwidth}{%
\raisebox{-0.9\height}{%
\phantom{.}
\includegraphics[height=18mm]{../../KUT000/URKFEIlogo_v0.1.pdf}%
}%
}%
\parbox[t][18mm][c]{0.7\textwidth}{%
\raggedleft

\sffamily
\fontsize{16pt}{18pt}
\fontseries{sbc}
\selectfont

\noindent
\textcolor[rgb]{0.75, 0.75, 0.75}{\textls[25]{\oznacenieCelku}}
}%

\noindent
\parbox[t][16mm][b]{0.5\textwidth}{%
\raggedright

\color{Gray}
\sffamily

\fontsize{12pt}{12pt}
\selectfont
\mesiacRok

\fontsize{6pt}{10pt}
\selectfont
github.com/OkoliePracovnehoBodu/KUT

\fontsize{8pt}{10pt}
\selectfont
\authorslabel




}%
\parbox[t][16mm][b]{0.5\textwidth}{%
\raggedleft

\sffamily

\fontsize{6pt}{6pt}
\selectfont

\textcolor[rgb]{0.68, 0.68, 0.68}{\oznacenieVerzie}


\fontsize{14pt}{14pt}
\selectfont

\bfseries


\includegraphics[height=12pt]{../../KUT000/KUT_logo_v0.1.pdf}%
{%
\textls[-50]{\KUTporadoveCislo}
}%
}%

% -----------------------------------------------------------------------------




\vspace{6mm}

% ---------------------------------------------
\sffamily
\bfseries
\fontsize{18pt}{21pt}
\selectfont

\begin{flushleft}
    TS: Príklad datasetu\\pre prevodovú charakteristiku
\end{flushleft}

\bigskip

% -----------------------------------------------------------------------------
\normalsize
\normalfont
% -----------------------------------------------------------------------------

\lstset{style=mystyle}










\noindent
\lettrine[lines=1, nindent=1pt, loversize=0.0]{C}{ieľom} 
textu je sprostredkovanie a opis meraní a dát súvisiacich so statickými vlastnosťami laboratórneho dynamického systému TS. Predchádzajúce texty súvisiace so zariadením TS: \href{run:../../KUT017/TeX/KUT017.pdf}{\textsf{KUT017}}, \href{run:../../KUT015/TeX/KUT015.pdf}{\textsf{KUT015}} 







\section{Meranie prevodovej charakteristiky}

V kontexte statických vlastností systému má vo všeobecnosti význam hovoriť o~prevodovej charakteristike systému. Prevodová charakteristika je závislosť ustálených hodnôt výstupného signálu systému od ustálených hodnôt vstupného signálu systému.

Je zrejmé, že prevodová charakteristika sa týka systémov s prívlastkom statické, teda takých, ktoré nie sú astatické.

Prevodová charakteristika, niekde sa nazýva aj statická charakteristika, teda charakterizuje systém len v~ustálených stavoch. Neobsahuje informáciu o dynamike systému.



\subsection{Návrh merania}

Z opisu predmetného dynamického systému [\href{run:../../KUT015/TeX/KUT015.pdf}{\textsf{KUT015}}] vyplýva, že systém má dva výstupné signály, a dva vstupné signály teda vo všeobecnosti ide o MIMO systém. Systém je možné vnímať ako SISO systém, ak budeme uvažovať, že jeden zo vstupov slúži len na nastavenie prevádzkovej podmienky systému.

V nasledujúcom sa uvažuje, že vstupným signálom je ten, ktorý ovláda ohrev špirály (Vstup~1). Prevádzková podmienka systému je realizovaná ventilátorom (Vstup~2).

Systém má dva výstupy – signály zo senzora 1 a senzora 2. Všetky signály (vstupný, výstupné aj signál prevádzkovej podmienky) nadobúdajú hodnoty v rozsahu $0$ až $10$ pričom ide o napäťové signály vo voltoch [V].





\subsubsection{Voľba ustálených hodnôt vstupov}

O predmetnom systéme je známe, že výstupy systému sa ustália vždy, ak sú vstupný signál a prevádzková podmienka ustálené. Pre vyšetrovanie ustálených stavov bol využitý celý rozsah vstupného signálu a vybrané hodnoty reprezentujúce prevádzkové podmienky.

Návrh uvažuje ustálené hodnoty vstupného signálu uvedené v tabuľke~\ref{tab:ustalene_hodnoty_vstupneho_signalu} a zároveň vybrané ustálené hodnoty prevádzkových podmienok podľa tabuľky~\ref{tab:ustalene_hodnoty_podmienok}.

\begin{center}

\vspace{-10pt}    
    
\tabcaption{Ustálené hodnoty vstupného signálu [V]}
\label{tab:ustalene_hodnoty_vstupneho_signalu}

\lstyle

\begin{tabular*}{\textwidth}{@{ \extracolsep{\fill}} ccccccccccc}
\toprule
0 & 1 & 2 & 3 & 4 & 5 & 6 & 7 & 8 & 9 & 10 \\
\bottomrule
\end{tabular*}


\tabcaption{Vybrané ustálené hodnoty signálu o prevádzkových podmienkach [V]}
\label{tab:ustalene_hodnoty_podmienok}

\lstyle

\begin{tabular*}{\textwidth}{@{ \extracolsep{\fill}} ccccc}
\toprule
2 & 4 & 6 & 8 & 10 \\
\bottomrule
\end{tabular*}

\end{center}


\subsubsection{Voľba časového intervalu pre ustálenie výstupu systému}

Empirické skúsenosti so systémom ukazujú, že predmetný systém  má relatívne pomalú dynamiku. Vo všeobecnosti sa systém dokáže úplne ustáliť približne do $1$ až $3$ minút po zmene vstupného signálu, avšak v niektorých prípadoch nemusí byť ani čas $10$ minút postačujúci na úplné ustálenie.

Ako kompromis medzi presnosťou merania a celkovou dĺžkou experimentu bol zvolený čas ustálenia $5$ minút pre ustálené hodnoty vstupného signálu $1$ až $10$ [V]. Pre nulovú hodnotu vstupného signálu bol z dôvodu pomalšieho poklesu teploty zvolený čas ustálenia $10$ minút. Celková dĺžka jedného merania tak dosahuje jednu hodinu.



\subsubsection{Postup merania}

Vzhľadom na uvedené voľby ustálených hodnôt a časového intervalu návrh predpokladá nasledovný postup.

\begin{enumerate}[leftmargin=0pt, labelsep=4mm, itemsep=0pt]
    \item Nastavenie prevádzkových podmienok na hodnotu z tabuľky~\ref{tab:ustalene_hodnoty_podmienok}.
    \item Postupná zmena vstupného signálu na hodnoty z tabuľky~\ref{tab:ustalene_hodnoty_vstupneho_signalu} so zvoleným časovým intervalom. Takúto postupnú zmenu vyjadruje nasledujúca tabuľka \ref{tab:postupna_zmena_vstupneho_signalu}.
    
\end{enumerate}


\begin{center}

\vspace{-10pt}    

\tabcaption{Postupná zmena vstupného signálu}
\label{tab:postupna_zmena_vstupneho_signalu}

\lstyle

\begin{tabular*}{\textwidth}{@{ \extracolsep{\fill}} cc}
\toprule
Čas zmeny vstupného signálu [s] & Hodnota vstupného signál [V] \\
\midrule
0 & 0 \\
600 & 1 \\
900 & 2 \\
1200 & 3 \\
1500 & 4 \\
1800 & 5 \\
2100 & 6 \\
2400 & 7 \\
2700 & 8 \\
3000 & 9 \\
3300 & 10 \\
\bottomrule
\end{tabular*}

\end{center}

Celková dĺžka merania je teda $3300 + 300 = 3600$ sekúnd (1 hodina) a počas tejto doby sú prevádzkové podmienky konštantné.















\section{Získané dáta}


Na základe uvedeného návrhu a postupu merania boli vytvorené skripty a simulačné schémy v rámci prostredia MATLAB - Simulink. Ich opis a dokumentácia nech sú nad rámec tohto textu. Výsledný dataset sa nachádza v adresári \texttt{dataRepo}.

Vizualizácia týchto dát je realizovaná pomocou Python skriptu v Jupyter notebooku \texttt{PY/job\_makeFigs.ipynb}.

Získané dáta všetky prevádzkové podmienky, ktoré sú v~tabuľke~\ref{tab:ustalene_hodnoty_podmienok}, sú vizualizované na nasledujúcich obrázkoch.















\paragraph{Meranie ustálených hodnôt pri prevádzkovej podmienke 2 [V]}

\begin{center}

    \vbox{%
        \makebox[\textwidth][c]{%
        \includegraphics{fj_01_exp1_panel_1.pdf}
        }

        \makebox[\textwidth][c]{%
        \includegraphics{fj_01_exp1_panel_3.pdf}
        }

        \makebox[\textwidth][c]{%
        \includegraphics{fj_01_exp1_panel_2.pdf}
        }

        \figcaption{ 
           Meranie pri nastavenej prevádzkovej podmienke $2$ [V].
           Dáta na obrázku dostupné ako dátové súbory v~adresári \texttt{PY/fig/fj\_01\_exp1\_panel\_\#.csv}
        }
        \label{fj_01_exp1}
    }%vbox

\end{center}

\vfill

\paragraph{Meranie ustálených hodnôt pri prevádzkovej podmienke 4 [V]}

\begin{center}

    \vbox{%
        \makebox[\textwidth][c]{%
        \includegraphics{fj_01_exp2_panel_1.pdf}
        }

        \makebox[\textwidth][c]{%
        \includegraphics{fj_01_exp2_panel_3.pdf}
        }

        \makebox[\textwidth][c]{%
        \includegraphics{fj_01_exp2_panel_2.pdf}
        }

        \figcaption{ 
            Meranie pri nastavenej prevádzkovej podmienke $4$ [V].
            Dáta na obrázku dostupné ako dátové súbory v~adresári \texttt{PY/fig/fj\_01\_exp2\_panel\_\#.csv}
        }

        \label{fj_01_exp2}
    }%vbox

\end{center}


\pagebreak














\paragraph{Meranie ustálených hodnôt pri prevádzkovej podmienke 6 [V]}

\begin{center}

    \vbox{%
        \makebox[\textwidth][c]{%
        \includegraphics{fj_01_exp3_panel_1.pdf}
        }

        \makebox[\textwidth][c]{%
        \includegraphics{fj_01_exp3_panel_3.pdf}
        }

        \makebox[\textwidth][c]{%
        \includegraphics{fj_01_exp3_panel_2.pdf}
        }

        \figcaption{ 
           Meranie pri nastavenej prevádzkovej podmienke $6$ [V].
           Dáta na obrázku dostupné ako dátové súbory v~adresári \texttt{PY/fig/fj\_01\_exp3\_panel\_\#.csv}
        }
        \label{fj_01_exp3}
    }%vbox

\end{center}

\vfill

\paragraph{Meranie ustálených hodnôt pri prevádzkovej podmienke 8 [V]}

\begin{center}

    \vbox{%
        \makebox[\textwidth][c]{%
        \includegraphics{fj_01_exp4_panel_1.pdf}
        }

        \makebox[\textwidth][c]{%
        \includegraphics{fj_01_exp4_panel_3.pdf}
        }

        \makebox[\textwidth][c]{%
        \includegraphics{fj_01_exp4_panel_2.pdf}
        }

        \figcaption{ 
            Meranie pri nastavenej prevádzkovej podmienke $8$ [V].
            Dáta na obrázku dostupné ako dátové súbory v~adresári \texttt{PY/fig/fj\_01\_exp4\_panel\_\#.csv}
        }

        \label{fj_01_exp4}
    }%vbox

\end{center}


\pagebreak











\paragraph{Meranie ustálených hodnôt pri prevádzkovej podmienke 10 [V]}

\begin{center}

    \vbox{%
        \makebox[\textwidth][c]{%
        \includegraphics{fj_01_exp5_panel_1.pdf}
        }

        \makebox[\textwidth][c]{%
        \includegraphics{fj_01_exp5_panel_3.pdf}
        }

        \makebox[\textwidth][c]{%
        \includegraphics{fj_01_exp5_panel_2.pdf}
        }

        \figcaption{ 
           Meranie pri nastavenej prevádzkovej podmienke $10$ [V].
           Dáta na obrázku dostupné ako dátové súbory v~adresári \texttt{PY/fig/fj\_01\_exp5\_panel\_\#.csv}
        }
        \label{fj_01_exp5}
    }%vbox

\end{center}

\vfill







% -----------------------------------------------------------------------------

\end{document}

% -----------------------------------------------------------------------------